\chapter{End-to-End Evaluation: Does Routing Pay?}
\label{ch:e2e}

\chapterhandoff{Chapter~\ref{ch:pred} showed the routing target is weakly
predictable and that most of the apparent predictability came from a benchmark
annotation. A classifier score is not a system result. This chapter prices the
policy against the alternatives a deployment actually has, and finds that the
choice of comparator explains more of the answer than the choice of learner
does.}

\section{The comparator ladder}
\label{sec:comparators}

A routed policy has to beat something, and which something is chosen determines
what the result means. We evaluate against four references, ordered by how hard
they are to beat.

\begin{enumerate}[leftmargin=1.5em]
\item \textbf{Best single mode.} The strongest fixed policy, selected with
      hindsight over the cell. Highest success, and among the most expensive.
\item \textbf{Always-cheapest.} Send every task to the cheapest mode. This
      requires no model, no features and no training: a deployment has it for free. It is also, by construction, a \emph{cost lower bound} over fixed
      policies.
\item \textbf{Oracle triage.} Retrospectively route hopeless tasks to the
      cheapest mode and everything else to the best mode. Not deployable;
      included as the ceiling of the half of the problem we can actually attempt.
\item \textbf{Learned triage.} The nested-CV router of
      Section~\ref{sec:protocol}, at the threshold its own inner folds selected.
\end{enumerate}

\begin{scopebox}
\textbf{What ``end to end'' means here, precisely.} The routed policies are
evaluated by \emph{replay}: for each task, the router's decision selects which
already-executed condition's outcome and cost are credited. We do not re-execute
the agent under a live router. This is a real limitation, and its direction is worth stating: replay omits any online routing overhead and any interaction
between routing and episode dynamics, so it is \emph{favourable} to the router.
A negative result under replay is therefore stronger than the same result under
live execution would be. Section~\ref{sec:overhead} bounds the omitted term.
\end{scopebox}

\section{The primary result}
\label{sec:primary-result}

\begin{table}[t]
\centering\small
\caption[End-to-end policy comparison, all eight cells]{\textbf{Every policy in
every cell, against the free fixed alternative.} $\Delta$ columns are against
always-cheapest. ``Learned (nested)'' is the deployment-faithful row. ``Learned
(swept)'' is the most favourable learned operating point available: its
threshold is chosen post hoc \emph{and} its best/cheapest modes are selected
from whole-cell outcomes, so it is optimistic on two axes at once. It is shown
so the conclusion cannot be blamed on either selection step. A cell
is a Pareto win only if $\Delta$SR $>0$ \emph{and} $\Delta$cost $<0$.}
\label{tab:e2e}
\begin{tabular}{@{}llrrrr@{}}
\toprule
Cell & policy & SR \% & cost & $\Delta$SR & $\Delta$cost \\
\midrule
\multirow{4}{*}{\cell{Classifieds$\cdot$235B}}
 & always-cheapest (\mVis{})   & 25.00 & 0.06481 & ref. & ref. \\
 & oracle triage               & 27.23 & 0.06312 & $+2.23$\pp & $-2.6$\% \\
 & learned (nested)            & 26.79 & 0.07129 & $+1.79$\pp & $+10.0$\% \\
 & learned (swept)  & 27.23 & 0.06730 & $+2.23$\pp & $+3.8$\% \\
\midrule
\multirow{4}{*}{\cell{Reddit$\cdot$235B}}
 & always-cheapest (\mVis{})   & \phantom{0}7.39 & 0.09807 & ref. & ref. \\
 & oracle triage               & 14.78 & 0.09998 & $+7.39$\pp & $+1.9$\% \\
 & learned (nested)            & 12.81 & 0.09928 & $+5.42$\pp & $+1.2$\% \\
 & learned (swept)  & 14.78 & 0.11045 & $+7.39$\pp & $+12.6$\% \\
\midrule
\multirow{4}{*}{\cell{Classifieds$\cdot$Qwen-4B}}
 & always-cheapest (\mVis{})   & 12.50 & 0.04316 & ref. & ref. \\
 & oracle triage               & 14.29 & 0.04858 & $+1.79$\pp & $+12.6$\% \\
 & learned (nested)            & 14.29 & 0.05867 & $+1.79$\pp & $+35.9$\% \\
 & learned (swept)  & 14.29 & 0.05874 & $+1.79$\pp & $+36.1$\% \\
\midrule
\multirow{4}{*}{\cell{Reddit$\cdot$Qwen-4B}}
 & always-cheapest (\mVis{})   & \phantom{0}2.46 & 0.05240 & ref. & ref. \\
 & oracle triage               & \phantom{0}7.39 & 0.05554 & $+4.93$\pp & $+6.0$\% \\
 & learned (nested)            & \phantom{0}6.40 & 0.07021 & $+3.94$\pp & $+34.0$\% \\
 & learned (swept)  & \phantom{0}7.39 & 0.07057 & $+4.93$\pp & $+34.7$\% \\
\midrule
\multirow{4}{*}{\cell{Classifieds$\cdot$Gemma-4B}}
 & always-cheapest (\mVis{})   & \phantom{0}2.23 & 0.07065 & ref. & ref. \\
 & oracle triage               & \phantom{0}2.23 & 0.07145 & $+0.00$\pp & $+1.1$\% \\
 & learned (nested)            & \phantom{0}1.34 & 0.07385 & $-0.89$\pp & $+4.5$\% \\
 & learned (swept)  & \phantom{0}2.23 & 0.07220 & $+0.00$\pp & $+2.2$\% \\
\midrule
\multirow{4}{*}{\cell{Reddit$\cdot$Gemma-4B}}
 & always-cheapest (\mVis{})   & \phantom{0}1.97 & 0.06833 & ref. & ref. \\
 & oracle triage               & \phantom{0}3.94 & 0.06974 & $+1.97$\pp & $+2.1$\% \\
 & learned (nested)            & \phantom{0}4.43 & 0.06910 & $+2.46$\pp & $+1.1$\% \\
 & learned (swept)  & \phantom{0}4.43 & 0.06950 & $+2.46$\pp & $+1.7$\% \\
\midrule
\multirow{4}{*}{\cell{Reddit (WA)$\cdot$235B}}
 & always-cheapest (\mDOM{})   & 26.92 & 0.07531 & ref. & ref. \\
 & oracle triage               & 35.58 & 0.07573 & $+8.66$\pp & $+0.6$\% \\
 & learned (nested)            & 32.69 & 0.08649 & $+5.77$\pp & $+14.8$\% \\
 & learned (swept)  & 35.58 & 0.08103 & $+8.66$\pp & $+7.6$\% \\
\midrule
\multirow{4}{*}{\cell{Reddit (WA)$\cdot$Qwen-4B}}
 & always-cheapest (\mVis{})   & \phantom{0}9.62 & 0.04468 & ref. & ref. \\
 & oracle triage               & 16.35 & 0.04814 & $+6.73$\pp & $+7.7$\% \\
 & learned (nested)            & 13.46 & 0.06428 & $+3.84$\pp & $+43.9$\% \\
 & learned (swept)  & 16.35 & 0.06331 & $+6.73$\pp & $+41.7$\% \\
\bottomrule
\end{tabular}
\end{table}

\begin{scopebox}
\textbf{Primary result.} In \textbf{0 of 8 cells} does the learned triage
router achieve both higher task success \emph{and} lower cost than the trivial
always-cheapest fixed policy. That joint condition --- better on both counts,
worse on neither --- is what we mean by a Pareto win throughout. The result
holds across two benchmarks, three models and two model families. This holds at the nested threshold and
at the most favourable in-sample threshold alike, so it is not an artefact of
operating-point selection.
\end{scopebox}

Every learned row in Table~\ref{tab:e2e} buys success with cost. The exchange
rate varies by more than an order of magnitude: \cell{red$\cdot$B2} pays 1.7\%
more for $+2.46\pp$, while \cell{wa\_red$\cdot$B1} pays 41.7\% more for
$+6.73\pp$ and \cell{cls$\cdot$B1} pays 36.1\% more for $+1.79\pp$. These are
trade-off points. Some of them a deployment might reasonably choose, but only after stating a price for a success, which our accounting does not do and which
the routing literature generally leaves implicit.

\section{The comparator explains more than the learner does}
\label{sec:comparator-effect}

The natural reading of ``0 of 8'' is that the learner was inadequate. That
reading is wrong, and the evidence against it is in the same table.

Normalise each cell to its always-cheapest policy and plot success against cost
(Figure~\ref{fig:f13}). The Pareto winning region is the upper-left quadrant:
more success \emph{and} less cost. \textbf{The retrospective oracle triage policy
enters that quadrant in only 1 of 8 cells} (\cell{cls$\cdot$B0}: $+2.23\pp$ at
$-2.6\%$). In the other seven the oracle, which knows the answer, also pays more than always-cheapest.

\begin{figure}[t]
  \centering
  \includegraphics[width=\textwidth]{fig_f13_dominance_plane.pdf}
  \caption[The dominance plane, normalised to always-cheapest]{%
  \textbf{Why ``0 of 8'' is a statement about the trade-off, not about the
  learner.} Each cell is normalised to its own always-cheapest policy at the
  origin; the shaded upper-left quadrant is the Pareto winning region. Learned
  policies sit right of the vertical axis: they cost more. So does the
  retrospective oracle, in seven of eight cells. Always-cheapest is the
  cheapest fixed policy available, so any chooser that raises success
  necessarily also spends more; whether that exchange is worth making requires
  a rate this dissertation does not assume, which is why the test is the joint
  one. The negative result is about the economics of this boundary,
  not about how good the router is at predicting.}
  \label{fig:f13}
\end{figure}

The explanation is structural, and it is worth stating precisely because the
obvious version of it is wrong. Always-cheapest minimises cost among fixed
policies, but only \emph{in expectation over the cell}: it is the mode with the lowest \emph{mean} episode cost, not the cheapest mode on every individual
task. A per-task router therefore \emph{can} undercut it, by sending a task to a
mode that is both more successful and, on that task, cheaper. This is
not hypothetical: on \cell{cls$\cdot$B0} the oracle does exactly that, gaining
$+2.23$\pp{} at $-2.6$\% cost.

So the winning region is narrow, not empty by construction. Entering it requires
the cheap mode to be \emph{wasteful} on a substantial subset of tasks, burning step budget on failures that a better mode would have resolved or abandoned
sooner. That happens, but on our data it happens often enough to pay in only one
cell of eight, and only for a policy that already knows the answers.

\begin{scopebox}
\textbf{Calibrated statement of the negative result.} Under this cost boundary,
on these benchmarks, at 2--36\% base success, \textbf{no profitable deployable operating point exists}, not because the router cannot be learned, but
because the region in which a router could win is nearly empty even for a policy
that knows the answers. This is a weaker claim than ``routing does not work''
and a more useful one: it identifies the condition (a cost-minimising fixed
comparator, in a regime where most tasks are hopeless) under which the enterprise
cannot pay.
\end{scopebox}

\section{Router overhead}
\label{sec:overhead}

Equation~\ref{eq:netsaving} requires the router's own cost. Our routed policies
are evaluated by replay, so this term is not in Table~\ref{tab:e2e}; we bound it
here to show it does not change the conclusion.

\begin{itemize}[leftmargin=1.5em]
\item \textbf{Feature extraction.} All 18 features are computed from the task
      instruction and the cheap structured observation by string operations and
      regular expressions. Both are already constructed for the cheap mode, so
      the marginal cost is a small constant per decision.
\item \textbf{Router inference.} A logistic regression over 18 standardised
      features: a dot product. At the amortised device rate used in our
      configuration ($5.6\times10^{-8}$\,\$/ms), this is several orders of
      magnitude below the $10^{-3}$\,\$ scale of a single agent step.
\item \textbf{Additional observation construction.} Zero for triage: the
      decision selects \emph{which} mode to construct, and exactly one is
      constructed either way. This is the property that makes triage the cheap
      half; a cascade that tries a mode and then escalates would not have it,
      and we measure that separately below.
\item \textbf{Error cost.} This is the term that dominates, and it is already
      inside Table~\ref{tab:e2e}. A false negative (a solvable task sent to the
      cheap mode) costs a success; a false positive (a hopeless task sent to the
      expensive mode) costs the price difference for a whole episode.
\end{itemize}

The three explicit overhead terms together are bounded by roughly $10^{-6}$\,\$
per decision against per-episode costs of $10^{-3}$ to $10^{-1}$\,\$. Adding
them moves no cell across the Pareto boundary: every learned row is already
1.1--43.9\% \emph{above} always-cheapest, and the overhead is four to five
orders of magnitude smaller than those gaps. The negative result does not depend
on the overhead accounting; it survives even the maximally generous version in
which routing is free.

\subsection{The cascade is the expensive alternative, and we priced it}

One obvious alternative to prediction is not to predict: run the cheapest mode,
detect failure, escalate. We evaluated a cost-ordered cascade with a perfect success detector, a strictly optimistic assumption, since real failure
detection on these tasks is itself unsolved.

It reaches the full oracle success in every cell, and costs $+315.9\%$ to
$+462.9\%$ against the best single mode. The reason is visible in the attempt
histograms: on \cell{cls$\cdot$B0}, 133 of 224 tasks reach all six attempts,
because no mode solves them and the cascade discovers this only after paying for
all six. In a regime where 48--93\% of tasks are hopeless, try-then-escalate
spends most of its budget confirming hopelessness.

This is worth stating positively as well as negatively: \textbf{the value of
prediction in this setting is exactly the value of not paying for the cascade},
and the cascade's price ($4$--$5\times$) is the size of the prize that
Chapter~\ref{ch:pred} could not claim.

\section{Site and backbone heterogeneity}
\label{sec:heterogeneity}

The aggregate is not driven by one cell. The 0-of-8 result reproduces on both benchmarks, on both VisualWebArena sites, and on all three backbones, and the
1-of-8 oracle result is likewise not concentrated in a single family: the one cell where the oracle wins (\cell{cls$\cdot$B0}) is the cell with the highest
solvable rate on VisualWebArena and the most minted labels, which is where
one would expect routing to work if it worked anywhere.

The two WebArena cells behave the same way despite a different corpus, a
different site, and a mode ordering in which an image-free mode is the strongest
single policy. That is the external-validation result: the failure transfers.

\begin{scopebox}
\textbf{``8 of 8'' is not eight independent trials, and we do not use it as
such.} The three backbones run the \emph{same} task sets, so within a benchmark
the cells are correlated through their shared tasks, and if base success rate is the mechanism, as Chapter~\ref{ch:disc} argues, then B1 and B2 are weaker
draws from the same regime rather than independent replications of it. The count
is a statement of consistency, not a sample size, and the effective number of
independent settings is closer to the two benchmarks than to eight cells. It is
reported because a result that reversed on some cells would be visible in it.
\end{scopebox}

\section{Robustness}
\label{sec:robustness}

\paragraph{Threshold.} Reported at both the nested threshold and the most
favourable in-sample threshold; the verdict is identical
(Table~\ref{tab:e2e}).

\paragraph{Label definition.} The which-mode supply result of
Table~\ref{tab:supply} was recomputed under an alternative label
(cheapest-successful by \emph{measured} episode cost rather than by prior cost
ordering). Under it, 5 of 6 cells fall below that minimum rather than 4 of 6. The conclusion does not turn on the definition, and the alternative is worse.

\paragraph{Cost metric.} The comparison is repeated under a cost estimand whose
\emph{meaning} differs across backbones: B0's per-1k schedule is a vendor price,
B1/B2's is a device-amortisation rate (Section~\ref{sec:metrics}). A vendor
price can weight images arbitrarily relative to text, so a cheapest-mode ranking
is not guaranteed to be stable across the two. On our data it is: \mVis{} is the
cheapest mode in seven of eight cells under the token schedule, and the one
exception (\cell{wa\_red$\cdot$B0}, where \mDOM{} is cheapest) is a WebArena
cell, not an API/local boundary effect. The 0-of-8 verdict holds separately
within the API cells (3 of 3) and within the local cells (5 of 5), so it does not
depend on comparing across the boundary.

\paragraph{Permutation budget.} $B$ was raised from 200 to 10{,}000 during this
work. At $B=200$ the reportable floor sat inside Holm's threshold, so the budget
could in principle have decided a verdict. At $B=10{,}000$ the smallest
reportable $p$ is $1.0\times10^{-4}$ against a tightest threshold of
$6.25\times10^{-3}$, so it cannot.

\paragraph{Noise.} The measured 12--14\% run-to-run discordance of
Section~\ref{sec:noise-floor} is a threat to positive claims, not to this one.
Noise degrades whatever signal exists; a negative result obtained despite it is
if anything understated.

\paragraph{What would overturn this.} A cell in which always-cheapest is \emph{not} close to cost-optimal, that is, where the cheap mode wastes steps on a substantial fraction of tasks, would open the winning region. Our
\cell{cls$\cdot$B0} is the one cell where this happens, and it is the one cell
where the oracle wins. That is the mechanism, and it is testable.

\section{What this chapter establishes}
\label{sec:ch6-summary}

At no threshold, in any of the eight cells, does the learned triage router Pareto-beat the free fixed alternative, and that verdict comes from a replay evaluation which omits the router's own overhead in its favour. Neither adding the overhead nor switching to the alternative label definition changes that.

More importantly, neither does replacing the learner with an oracle. In seven of
eight cells a policy that knows every outcome still cannot beat always-cheapest
on both axes. The negative result is therefore about the economics of the boundary, a cost-minimising comparator in a regime where most tasks are unsolvable, and not about predictive capacity.

\chapterhandoff{The end-to-end evaluation shows that the opportunity identified
in Chapter~\ref{ch:repr} does not convert into deployable benefit, and locates
the reason in the comparator rather than the learner. What remains is to explain
why the predictive half failed on its own terms, to say how far these
conclusions travel, and to state precisely which regime they describe.
Chapter~\ref{ch:disc} takes up all three, and closes the dissertation.}
